\subsection{Infrastructure for RL Training of LLMs}
\label{sec:infrastructure-rl-training}

Reinforcement learning (RL) for large language models introduces infrastructure demands that go substantially beyond those of supervised fine-tuning. A single RL training step---whether PPO, GRPO, or REINFORCE++---requires (i) generating multiple candidate completions from the current policy (on-policy sampling), (ii) scoring those completions with reward models or verifiers, and (iii) computing policy-gradient updates against the scored rollouts. Each phase imposes distinct computational requirements: sampling is autoregressive and memory-bound, reward evaluation may involve external tools or secondary models, and gradient computation is compute-bound and must operate over the full model. Crucially, all three phases must share or synchronize model weights, and for on-policy methods such as GRPO the sampling weights must remain close to the training weights to avoid excessive staleness. This tight coupling between generation and training, combined with the need to hold multiple model replicas in GPU memory, makes RL training of LLMs one of the most infrastructure-intensive workloads in modern machine learning. The tools and frameworks surveyed in this section address these challenges at different layers of the stack.

\subsubsection{TRL: Transformer Reinforcement Learning}

TRL (Transformer Reinforcement Learning) is an open-source library developed by Hugging Face that provides a unified interface for post-training language models with reinforcement learning methods \cite{vonwerra2022trl}. The library is tightly integrated with the \texttt{transformers} and \texttt{accelerate} ecosystems, enabling researchers to apply RL algorithms to any Hugging Face--compatible model with minimal boilerplate. TRL implements a comprehensive suite of post-training algorithms including Proximal Policy Optimization (PPO), Direct Preference Optimization (DPO), Group Relative Policy Optimization (GRPO), Reward Modeling, and Supervised Fine-Tuning (SFT).

Of particular relevance to this work is the \texttt{GRPOTrainer} class, which implements the GRPO algorithm introduced in DeepSeekMath \cite{shao2024deepseekmath}. The \texttt{GRPOTrainer} manages the full RL loop: for each prompt in a batch, it generates a group of $G$ completions, computes per-completion rewards via user-defined reward functions, normalizes rewards within each group to obtain advantages, and performs a clipped policy-gradient update. TRL's modular design allows users to supply custom reward functions---ranging from rule-based verifiers to learned reward models---and to control generation parameters such as temperature, top-$k$, and sampling strategy. Distributed training is supported through DeepSpeed ZeRO integration and Fully Sharded Data Parallelism (FSDP), and TRL can offload generation to vLLM for significantly higher sampling throughput. The library has become the de facto starting point for academic GRPO experiments at moderate scale.

\subsubsection{OpenRLHF}

OpenRLHF is a high-performance, production-ready RLHF framework that combines Ray, vLLM, and DeepSpeed into a unified distributed architecture \cite{hu2024openrlhf}. Whereas TRL targets single-node or modest multi-GPU setups, OpenRLHF is explicitly designed for large-scale distributed RL training on clusters with tens to hundreds of GPUs. The framework introduces a unified agent-based design paradigm in which the RL loop is decomposed into independently schedulable actors---policy, reference model, reward model, and critic---each of which can be placed on separate GPU groups and scaled independently via Ray.

A key innovation in OpenRLHF is its tight integration with vLLM for on-policy generation. During the rollout phase, vLLM serves as the inference backend with PagedAttention and continuous batching, achieving high throughput even for long chain-of-thought completions. Model weights are synchronized between the training engine (which uses DeepSpeed ZeRO Stage 3 for memory-efficient gradient computation) and the vLLM inference engine via Ray's object store, enabling near-on-policy sampling without requiring the entire model to be duplicated in memory. OpenRLHF supports PPO, REINFORCE++, GRPO, RLOO, and DAPO algorithms, along with both single-turn and multi-turn agent-based RL. The framework's Hybrid Engine mode allows training and generation to share the same GPU pool by dynamically switching between training and inference modes, further reducing hardware requirements.

\subsubsection{DeepSpeed and FSDP}

The memory and compute requirements of RL training at scale necessitate distributed training frameworks that can efficiently partition model state across many GPUs. DeepSpeed, developed by Microsoft, introduces the Zero Redundancy Optimizer (ZeRO) \cite{rajbhandari2020zero}, which partitions optimizer states (Stage 1), gradients (Stage 2), and model parameters (Stage 3) across data-parallel ranks, reducing per-GPU memory consumption from $O(1)$ to $O(1/N)$ for $N$ GPUs while preserving the simplicity of data parallelism. ZeRO-Offload and ZeRO-Infinity further extend this by offloading state to CPU memory or NVMe storage, enabling training of models with hundreds of billions of parameters on commodity hardware. DeepSpeed also provides Automatic Tensor Parallelism (AutoTP), which combines tensor-parallel weight sharding within nodes with ZeRO-based data parallelism across nodes, achieving hybrid 3D parallelism without requiring manual model partitioning.

PyTorch's Fully Sharded Data Parallelism (FSDP) provides a native alternative to ZeRO within the PyTorch ecosystem. FSDP shards model parameters, gradients, and optimizer states across data-parallel workers, all-gathering parameters on demand for forward and backward passes. Both DeepSpeed and FSDP are critical enablers for RL training: during the gradient computation phase of GRPO, the full model must participate in forward and backward passes over potentially hundreds of rollouts per batch, and without efficient sharding the memory footprint would be prohibitive. TRL, OpenRLHF, and most other RL training frameworks build directly on one or both of these backends.

\subsubsection{vLLM: High-Throughput Inference for On-Policy Sampling}

On-policy RL algorithms require generating fresh completions from the current policy at every training step. For large models producing long chain-of-thought outputs, this generation phase can account for up to 90\% of total training wall-clock time, making inference throughput the primary bottleneck in the RL loop. vLLM \cite{kwon2023vllm} is an open-source inference engine that addresses this bottleneck through two core innovations: \textit{PagedAttention}, which manages the key-value (KV) cache using a paging mechanism analogous to virtual memory in operating systems, eliminating memory fragmentation and enabling near-optimal GPU memory utilization; and \textit{continuous batching}, which dynamically admits and retires requests from the batch as they complete, avoiding the throughput loss caused by padding in static batching.

Together, these techniques allow vLLM to achieve 10--24$\times$ higher throughput than naive Hugging Face Transformers inference, with support for tensor parallelism across multiple GPUs for serving models that exceed single-GPU memory. In the context of RL training, vLLM is deployed as the sampling backend: the training framework (e.g., OpenRLHF or TRL) dispatches generation requests to a vLLM instance, collects the completions and their log-probabilities, and then performs the policy-gradient update. Weight synchronization between the training engine and the vLLM server is handled either through shared memory, Ray's object store, or periodic checkpoint reloading, depending on the framework.

\subsubsection{Tinker: Training-as-a-Service for RL Research}

Tinker, developed by Thinking Machines Lab (the AI research startup founded by former OpenAI executive Mira Murati), represents a fundamentally different approach to RL training infrastructure \cite{tinker2025}. Rather than requiring researchers to provision and manage GPU clusters, Tinker exposes a cloud-based Python API that abstracts distributed training into four primitive operations: \texttt{forward\_backward} (compute gradients), \texttt{optim\_step} (update weights), \texttt{sample} (generate tokens), and \texttt{save\_state} (checkpoint). Researchers write training loops that execute locally on a CPU, while Tinker orchestrates the actual computation on managed distributed GPU clusters behind the API.

Tinker employs LoRA (Low-Rank Adaptation) for all fine-tuning, training a compact adapter rather than updating full model weights. The platform's research demonstrates that with proper configuration, LoRA matches the learning dynamics of full fine-tuning while requiring substantially less compute. Tinker supports a broad range of open-source models, from compact dense models such as Llama-3.2-1B to large Mixture-of-Experts architectures such as Qwen3.5-397B-A17B and DeepSeek-V3.1. Pricing follows a per-million-token model across prefill, sampling, and training operations. By handling scheduling, fault tolerance, multi-node orchestration, and resource management, Tinker enables researchers without access to large compute clusters to conduct RL experiments on frontier-scale models, significantly democratizing access to RL training research.

\subsubsection{Atropos and the Tinker--Atropos Integration}

Atropos, developed by Nous Research, is an environment microservice framework for asynchronous reinforcement learning with LLMs \cite{atropos2025}. Atropos decouples the RL loop into three independently deployable components: \textit{environments} (which generate prompts, parse completions, and compute rewards), a \textit{trajectory API} (which queues and batches scored rollouts), and a \textit{trainer} (which consumes batches and updates the policy). Environments are implemented as lightweight HTTP services following the \texttt{AtroposBaseEnv} pattern, supporting diverse task types including mathematical reasoning (GSM8K), code generation (MBPP, HumanEval), tool calling, RLAIF, and multi-turn interactions.

The \textit{tinker-atropos} integration layer bridges Atropos environments with the Tinker training API, enabling any Atropos environment to leverage Tinker's managed compute without modification. The integration exposes OpenAI-compatible inference endpoints that route sampling requests through Tinker's \texttt{SamplingClient}, while Atropos's \texttt{ManagedServer} pattern handles log-probability tracking across multi-turn interactions for importance-sampling loss computation. Experimental results on the GSM8k benchmark with Llama-3.1-8B-Instruct demonstrate that the Tinker--Atropos integration achieves comparable reward trajectories to the native Tinker training loop, with marginally lower per-step training time due to Atropos's asynchronous architecture enabling overlapped inference and gradient computation. The integration reaches a reward plateau of 0.9 in approximately 2 minutes compared to 3 minutes for the synchronous Tinker trainer, illustrating the throughput benefits of decoupled, asynchronous RL orchestration.

\subsubsection{Why Infrastructure Matters}

The infrastructure challenges of RL training for LLMs can be summarized along three axes. First, the \textit{generation--training loop} requires tight integration between fast autoregressive inference and gradient-based optimization, two workloads with fundamentally different computational profiles (memory-bound vs.\ compute-bound). Second, \textit{multi-model coordination} demands that policy models, reference models, reward models, and critics coexist in GPU memory or be efficiently swapped, often requiring 2--4$\times$ the memory of supervised fine-tuning. Third, \textit{distributed compute orchestration} must handle heterogeneous workloads (generation, scoring, training) with variable latencies across potentially hundreds of GPUs while maintaining weight freshness for on-policy guarantees.

The frameworks surveyed in this section address these challenges through complementary strategies: DeepSpeed and FSDP reduce per-GPU memory through state sharding; vLLM maximizes generation throughput through PagedAttention and continuous batching; OpenRLHF orchestrates the full RL loop across distributed Ray actors; TRL provides an accessible entry point with deep Hugging Face integration; and Tinker with Atropos abstracts the entire infrastructure layer into a managed API with decoupled environment services. Together, these tools constitute the infrastructure stack that makes GRPO training at scale practically feasible, and the choice among them reflects a trade-off between control, scalability, and ease of use that researchers must navigate based on their available compute and experimental requirements.
